% ===========================================================================
% 04 - Architecture Decision Records (ADR) im Code
% ===========================================================================

% --- ADR-01 -----------------------------------------------------------------
\begin{frame}[fragile]{ADR-01 -- Fachlicher Schnitt statt technischer Layer}
  \adrbadge{ADR-01}
  \begin{columns}[T,onlytextwidth]
    \begin{column}{0.5\textwidth}
      \textbf{Entscheidung}
      \begin{itemize}
        \item Services werden nach \textbf{Business Capabilities} geschnitten.
        \item Kein technischer Layer-Schnitt (kein ``Controller-Service'').
      \end{itemize}

      \vspace{0.25cm}
      \textbf{Nachweis: Modulstruktur}
      \begin{lstlisting}[style=bash]
services/
  configuration-service/          -> Engine-Konfiguration
  analysis-management-service/    -> AnalysisRun, Status, Resultate
  fluid-analysis-service/         -> Oil-/Fuel-Analyse
  thermal-analysis-service/       -> thermische Analyse
  electrical-analysis-service/    -> elektrische Analyse
  engine-management-analysis-service/ -> EMS-Analyse
      \end{lstlisting}
    \end{column}
    \begin{column}{0.46\textwidth}
      \textbf{Interne Schichtung (je Service)}
      \begin{lstlisting}[style=bash]
fluid/
  api/             AnalysisController, AnalysisCommand
  application/     AnalysisWorker
  infrastructure/  NextServiceClient,
                   AnalysisManagementClient
  FluidAnalysisServiceApplication.java
      \end{lstlisting}

      \vspace{0.2cm}
      \begin{block}{Wirkung}
        \small Vermeidung des Anti-Patterns \textbf{Wrong Cut} (QR-06).
      \end{block}
    \end{column}
  \end{columns}
\end{frame}

% --- ADR-02 -----------------------------------------------------------------
\begin{frame}[fragile]{ADR-02 -- REST für die Kommunikation}
  \adrbadge{ADR-02}
  \begin{columns}[T,onlytextwidth]
    \begin{column}{0.52\textwidth}
      \begin{lstlisting}[style=java]
@RestController
@RequestMapping("/internal/analyses")
public class AnalysisController {

    private final AnalysisWorker worker;

    @PostMapping
    public ResponseEntity<Void> start(
            @RequestBody AnalysisCommand command) {
        worker.execute(command);
        return ResponseEntity.accepted().build();
    }
}
      \end{lstlisting}
      \srcfile{services/fluid-analysis-service/.../api/AnalysisController.java}
    \end{column}
    \begin{column}{0.44\textwidth}
      \textbf{Erläuterung}
      \begin{itemize}
        \item Alle Services kommunizieren über HTTP/REST mit definierten DTOs.
        \item Kein Service greift auf interne Klassen eines anderen zu.
        \item Ausgehende Aufrufe über \code{RestClient}.
      \end{itemize}
      \vspace{0.2cm}
      \begin{block}{Nachweis}
        \small \code{*Controller} (eingehend) + \code{RestClient}-Clients (ausgehend) in jedem Service.
      \end{block}
    \end{column}
  \end{columns}
\end{frame}

% --- ADR-03 -----------------------------------------------------------------
\begin{frame}[fragile]{ADR-03 -- Choreographie statt Orchestration}
  \adrbadge{ADR-03}
  \begin{columns}[T,onlytextwidth]
    \begin{column}{0.54\textwidth}
      \textbf{Start nur des Anchor-Algorithmus}
      \begin{lstlisting}[style=java]
public AnalysisRun start(String configurationId) {
    ConfigurationSnapshot configuration =
        configurationClient.get(configurationId);
    AnalysisRun run = AnalysisRun.start(
        "A-" + UUID.randomUUID(), configurationId);
    AlgorithmExecution fluid = run.execution(AlgorithmName.FLUID);
    fluid.updateStatus(AnalysisStatus.RUNNING, null);

    // Commit before external HTTP call (avoids callback race).
    run = repository.saveAndFlush(run);

    serviceStarter.start(AlgorithmName.FLUID,
        new AnalysisCommand(run.getAnalysisId(),
                            configuration, List.of()));
    return run;
}
      \end{lstlisting}
      \srcfile{.../application/AnalysisApplicationService.java}
    \end{column}
    \begin{column}{0.42\textwidth}
      \textbf{Weitergabe durch den Service}
      \begin{lstlisting}[style=java]
@Async
public void execute(AnalysisCommand command) {
    managementClient.reportStatus(
        command.analysisId(), ALGORITHM, "RUNNING", null);
    if (!pause(command.analysisId())) return;

    boolean ok = valid(command.configuration()
                       .get("oilSystem"))
              && valid(command.configuration()
                       .get("fuelSystem"));
    String result = ok ? "OK" : "FAILED";
    managementClient.reportResult(command.analysisId(),
        ALGORITHM, ok ? "READY" : "FAILED", result, null);

    if (ok) {
        nextServiceClient.startNext(command, result);
    }
}
      \end{lstlisting}
      \srcfile{.../application/AnalysisWorker.java}
    \end{column}
  \end{columns}
\end{frame}

% --- ADR-04 -----------------------------------------------------------------
\begin{frame}[fragile,shrink=5]{ADR-04 -- Circuit Breaker mit Resilience4j}
  \adrbadge{ADR-04}
  \begin{columns}[T,onlytextwidth]
    \begin{column}{0.54\textwidth}
      \begin{lstlisting}[style=java]
@CircuitBreaker(name = "analysisServiceStarter")
public void start(AlgorithmName algorithm,
                  AnalysisCommand command) {
    String baseUrl = switch (algorithm) {
        case FLUID -> fluidUrl;
        case THERMAL -> thermalUrl;
        case ELECTRICAL -> electricalUrl;
        case ENGINE_MANAGEMENT -> engineManagementUrl;
    };

    builder.baseUrl(baseUrl).build()
            .post().uri("/internal/analyses")
            .body(command).retrieve().toBodilessEntity();
}
      \end{lstlisting}
      \srcfile{.../infrastructure/AnalysisServiceStarter.java}
    \end{column}
    \begin{column}{0.42\textwidth}
      \textbf{Erläuterung}
      \begin{itemize}
        \item Umsetzung von \textbf{Isolation of Failures}.
        \item Der Fehler bleibt lokal im aufrufenden Service behandelbar.
        \item Zwei Breaker-Instanzen:
          \begin{itemize}
            \item \code{analysisServiceStarter} (Management $\rightarrow$ Analyse)
            \item \code{nextService} (Analyse $\rightarrow$ nächster Service)
          \end{itemize}
      \end{itemize}
      \begin{block}{Detail}
        \small Vollständiger Code + Fallback: siehe TA-04.
      \end{block}
    \end{column}
  \end{columns}
\end{frame}

% --- ADR-05 -----------------------------------------------------------------
\begin{frame}{ADR-05 -- Container pro Microservice}
  \adrbadge{ADR-05}
  \begin{columns}[T,onlytextwidth]
    \begin{column}{0.5\textwidth}
      \textbf{Entscheidung}
      \begin{itemize}
        \item Jeder Service besitzt ein eigenes \code{Dockerfile}.
        \item Jeder Service läuft in einem eigenen Container.
      \end{itemize}

      \vspace{0.2cm}
      \textbf{Nachweis}
      \begin{itemize}
        \item 6 $\times$ \code{Dockerfile} (Multi-Stage).
        \item \code{docker-compose.yml} mit 6 Services und Ports 8081--8086.
      \end{itemize}
    \end{column}
    \begin{column}{0.46\textwidth}
      \begin{center}
      \begin{tikzpicture}[
        every node/.style={font=\scriptsize},
        box/.style={draw=darkblue,rounded corners,fill=lightblue,minimum width=42mm,minimum height=6mm,align=center}
      ]
      \node[box] (a) {configuration-service :8081};
      \node[box,below=1.5mm of a] (b) {analysis-management :8082};
      \node[box,below=1.5mm of b] (c) {fluid-analysis :8083};
      \node[box,below=1.5mm of c] (d) {thermal-analysis :8084};
      \node[box,below=1.5mm of d] (e) {electrical-analysis :8085};
      \node[box,below=1.5mm of e] (f) {engine-management :8086};
      \end{tikzpicture}
      \end{center}

      \vspace{0.1cm}
      \begin{block}{Wirkung}
        \small Unterstützt Independent Deployability (QR-03).
      \end{block}
    \end{column}
  \end{columns}
\end{frame}

% --- ADR-06 -----------------------------------------------------------------
\begin{frame}[fragile,shrink=3]{ADR-06 -- Keine Shared Persistence}
  \adrbadge{ADR-06}
  \begin{columns}[T,onlytextwidth]
    \begin{column}{0.5\textwidth}
      \textbf{Getrennte Entities}
      \begin{lstlisting}[style=java]
@Entity
@Table(name = "engine_configurations")
public class EngineConfiguration {
    @Id
    private String configurationId;
    private String oilSystem;
    private String fuelSystem;
    // ...
}
      \end{lstlisting}
      \srcfile{configuration-service/.../domain/EngineConfiguration.java}

      \begin{lstlisting}[style=java]
@Entity
@Table(name = "analysis_runs")
public class AnalysisRun {
    @Id
    private String analysisId;
    // ...
}
      \end{lstlisting}
      \srcfile{analysis-management-service/.../domain/AnalysisRun.java}
    \end{column}
    \begin{column}{0.46\textwidth}
      \textbf{Integration über DTO}
      \begin{lstlisting}[style=java]
public record ConfigurationSnapshot(
        String configurationId,
        String oilSystem,
        String fuelSystem,
        String coolingSystem,
        String electricalSystem,
        String engineManagementSystem) {
}
      \end{lstlisting}
      \srcfile{.../infrastructure/ConfigurationSnapshot.java}
      \begin{block}{Nachweis stateless}
        \small Analyse-Services haben \textbf{keinen} \code{spring.datasource}-Block in \code{application.yml}.
      \end{block}
    \end{column}
  \end{columns}
\end{frame}
